\documentclass[11pt,a4paper]{article}
\usepackage[T1]{fontenc}
\usepackage{lmodern,amsmath,amssymb,amsthm,mathtools,hyperref}
\usepackage[margin=27mm]{geometry}
\setlength{\emergencystretch}{2em}
\newtheorem{theorem}{Theorem}[section]
\newtheorem{lemma}[theorem]{Lemma}
\newtheorem{corollary}[theorem]{Corollary}
\theoremstyle{definition}\newtheorem{definition}[theorem]{Definition}
\newcommand{\Q}{\mathbb Q}\newcommand{\Z}{\mathbb Z}
\title{A local obstruction to a quartic pole equation}
\author{}\date{}
\begin{document}\maketitle
\begin{abstract}
The equation $y^2=T^2+\kappa R(1+Rs^2)^2$ gives the strict
valuation inequality required to extract local tests from a reciprocal
trace. Nevertheless, with $\kappa=1$ its associated genus-one equation
has no point over $\Q_{1013}$ for an explicit infinite family of
prepared positive inputs tending to zero. The inputs have square
reduced denominators and satisfy all prescribed squareclass and
valuation conditions at $2,3,17$. Thus these preparation conditions
alone do not give completeness for this equation. More generally,
no finite set of fixed nonzero rational coefficients $\kappa$
is complete on all prepared inputs. The result does
not decide completeness on a subfamily selected together with the
comparison primes.
\end{abstract}

\section{The valuation implication}
We normalize the valuation $v_p$ of $\Q_p$ by $v_p(p)=1$ and put
$v_p(0)=+\infty$.

\begin{lemma}\label{lem:extraction}
Let $p$ be an odd prime and let $R,\kappa\in\Q_p^\times$ satisfy
$v_p(\kappa)=0$ and $\nu=v_p(R)\in2\Z+1$. If $y,s,T\in\Q_p$
satisfy
\begin{equation}\label{eq:candidate}
                   y^2=T^2+\kappa R(1+Rs^2)^2,
\end{equation}
then $T\ne0$ and
\begin{equation}\label{eq:strict}
                    2v_p(T)<\nu.
\end{equation}
\end{lemma}
\begin{proof}
If $s\ne0$, then $v_p(Rs^2)=\nu+2v_p(s)$ is odd and hence
nonzero. Therefore
\[
 u:=v_p(1+Rs^2)=\min(0,\nu+2v_p(s))\le0.
\]
For $s=0$ the same conclusion holds with $u=0$. In particular
$1+Rs^2\ne0$. The second term on the right of
\eqref{eq:candidate} has the odd valuation $\nu+2u$.
If $T=0$, this cannot be the valuation of a square. If $T\ne0$,
its square has even valuation and cannot cancel the other term.
The smaller valuation must be even, so
$2v_p(T)<\nu+2u\le\nu$.
\end{proof}

\begin{corollary}\label{cor:trace}
Let $K/\Q$ be quadratic and unramified at $p$, and let
$W\in K^\times$. Put
\[
 T=\operatorname{Tr}_{K/\Q}(W^{-1}),\qquad
 \lambda_p=\max_{\mathfrak p\mid p}v_{\mathfrak p}(W),
\]
where the local valuations have $v_{\mathfrak p}(p)=1$.
Under the hypotheses of Lemma~\ref{lem:extraction}, a rational
solution of~\eqref{eq:candidate} implies
\[
                         \lambda_p>-\nu/2.
\]
\end{corollary}
\begin{proof}
At a split prime the two terms of the trace have valuations at least
$-\lambda_p$. At an inert prime conjugation preserves the unique
valuation and gives the same conclusion. The nonarchimedean
triangle inequality yields $v_p(T)\ge-\lambda_p$.
Apply~\eqref{eq:strict}.
\end{proof}

Since $T\ne0$, the changes of variables
\[
             b=R/T^2,\qquad v=Ts,\qquad Y=y/T
\]
transform~\eqref{eq:candidate} exactly into
\begin{equation}\label{eq:quartic}
                   Y^2=1+\kappa b(1+bv^2)^2.
\end{equation}
For $b,\kappa>0$ the polynomial on the right is a separable
quartic: its leading and constant coefficients are nonzero, and
the discriminant of its quadratic expression in $v^2$ is
$-4\kappa b^3\ne0$. Its smooth projective double cover of
the $v$-line consequently has genus one.

\section{A local obstruction with even negative valuation}

\begin{lemma}\label{lem:negative}
Let $p\equiv1\pmod4$ be prime. Suppose that $\kappa$ is a
$p$-adic unit with square residue, and write
\[
                  b=p^{-2m}u,\qquad m\ge1,
                  \quad u\in\Z_p^\times,
\]
with $u$ having nonsquare residue. Then~\eqref{eq:quartic}
has no solution in $\Q_p$. Its smooth projective model also
has no $\Q_p$-point.
\end{lemma}
\begin{proof}
Since $v_p(\kappa b)<0$, one has $1+\kappa b\ne0$.
The discriminant $-4\kappa b^3$ is nonzero, so the quartic is
separable even when $\kappa$ is negative.
Consider $v\ne0$. If $v_p(v)<m$, then $v_p(bv^2)<0$ and
$v_p(1+bv^2)=v_p(bv^2)<0$. If $v_p(v)>m$, then
$v_p(1+bv^2)=0$. In the remaining case $v_p(v)=m$, the residue
of $bv^2$ is a nonsquare. Since $-1$ is a square modulo $p$,
the sum $1+bv^2$ is a unit. The case $v=0$ also gives a unit.
Thus in every case
\[
          v_p(1+bv^2)\le0,\qquad 1+bv^2\ne0.
\]
It follows that $\kappa b(1+bv^2)^2$ has negative even
valuation and nonsquare leading unit. Adding one preserves
both properties. The result cannot be a square in $\Q_p$.

At infinity, the two points of the smooth quartic are defined
over the extension obtained by adjoining a square root of the
leading coefficient $\kappa b^3$. Its valuation is even and
its leading unit is again a nonsquare. Neither point is defined
over $\Q_p$.
\end{proof}

\section{Prepared counterexamples}

\begin{definition}
A rational number $b$ is prepared if its reduced denominator is
a square and
\[
\begin{gathered}
 0<b<17,\qquad
 b\in\Q_2^{\times2}\cap\Q_3^{\times2}\cap\Q_{17}^{\times2},\\
 v_2(b)\ge3,\qquad v_3(b)\ge1,
                  \qquad v_{17}(b)\ge2.
\end{gathered}
\]
These are the preparation conditions of the pole-realization
input in~\cite[Section~3]{Dependencies}.
\end{definition}

\begin{theorem}\label{thm:family}
For every integer $m\ge1$, the number
\[
                         b_m=409\left(\frac{204}{1013^m}\right)^2
\]
is prepared, but the smooth projective curve
\begin{equation}\label{eq:family}
                         Y^2=1+b_m(1+b_mv^2)^2
\end{equation}
has no point over $\Q_{1013}$ and hence no rational point.
The numbers $b_m$ are distinct and tend to zero.
\end{theorem}
\begin{proof}
Trial division certifies that $1013$ is prime, and
$\gcd(409\cdot204,1013)=1$. Thus the reduced denominator of
$b_m$ is $1013^{2m}$. At $m=1$ one has
\[
 b_1=\frac{17020944}{1026169},\qquad
 17\cdot1026169-17020944=423929>0.
\]
Every $b_m$ lies in $(0,17)$, and
\[
          v_2(b_m)=4,\qquad v_3(b_m)=2,
                         \qquad v_{17}(b_m)=2.
\]
The congruences $409\equiv1\pmod8$, $409\equiv1\pmod3$
and $409\equiv1\pmod{17}$ show that $409$ is a square in
each of $\Q_2,\Q_3,\Q_{17}$. Its product with the displayed
rational square therefore satisfies all three local squareclass
conditions.

At $p=1013$ one has $p\equiv1\pmod4$ and
\[
                        409^{506}\equiv1012\pmod{1013}.
\]
The unit $409\cdot204^2$ is consequently a nonsquare modulo
$1013$, and $v_{1013}(b_m)=-2m$.
Lemma~\ref{lem:negative}, with $\kappa=1$, applies.
Finally $b_{m+1}=b_m/1013^2$, proving the last assertion.
\end{proof}

\begin{corollary}\label{cor:finite}
Let $C$ be a finite nonempty subset of $\Q^\times$. There are
arbitrarily small prepared $b$ for which none of the equations
\eqref{eq:quartic}, with $\kappa\in C$, has a rational point,
even on its smooth projective model.
\end{corollary}
\begin{proof}
Write $\kappa=a_\kappa/d_\kappa$ in reduced form with
$d_\kappa>0$, and put
\[
             M=408\prod_{\kappa\in C}|a_\kappa d_\kappa|.
\]
Choose a prime divisor $p$ of the cyclotomic value $\Phi_M(M)>1$.
It satisfies $p\nmid M$, since $\Phi_M(0)=1$. The cyclotomic
order lemma gives $p\equiv1\pmod M$: in characteristic $p$
the polynomial $x^M-1$ is separable, so its cyclotomic factors
are pairwise coprime and a root of $\Phi_M$ has exact order $M$.
In particular $p\equiv1\pmod8$ and $p\equiv1\pmod r$ for
every odd prime divisor $r$ of a numerator or denominator above.
Quadratic reciprocity shows that each such $r$, as well as
$-1$ and $2$, is a square modulo $p$. Thus every $\kappa\in C$
is a $p$-adic unit with square residue.

The Chinese remainder theorem supplies a positive integer $c$
which is congruent to one modulo $408$ and to a nonzero
nonsquare modulo $p$. For all sufficiently large $m$, the numbers
\[
                         b_m=c(204/p^m)^2
\]
are in $(0,17)$, have reduced denominator $p^{2m}$ and
valuations $4,2,2$ at $2,3,17$, and are squares in the three
corresponding completions. Hence they are prepared. At $p$
they have valuation $-2m$ and nonsquare leading unit
$c\cdot204^2$. Lemma~\ref{lem:negative} applies simultaneously
to every $\kappa\in C$, and $b_m\to0$.
\end{proof}

The valuation implication in Corollary~\ref{cor:trace} is valid
independently of this failure of completeness. Theorem~\ref{thm:family}
shows that the stated preparation conditions do not suffice to
replace the pole-realization theorem by~\eqref{eq:quartic} with
$\kappa=1$; Corollary~\ref{cor:finite} excludes any finite choice
of fixed nonzero rational coefficients. Neither result shows that
these counterexamples occur
among the inputs obtained from simultaneous comparison-prime
selection. No variable-count reduction follows without a realization
argument on the required selected inputs.

\section{Exact verification}
The standard-library Python program
\begin{center}\small
\nolinkurl{check_quartic_pole_obstruction.py}
\end{center}
certifies primality, the Euler congruence, the preparation congruences,
the exact first input, and the valuation case distinctions. It also
checks every residue modulo $1013$ in the potentially cancelling
case and saves its output in
\begin{center}\small
\nolinkurl{quartic-pole-obstruction.json}.
\end{center}
The proofs above establish the assertions for every $p$-adic
argument and every $m\ge1$; the certificate records the finite
arithmetic used in those proofs. No rank computation or conjectural
input is needed for the obstruction.

\begin{thebibliography}{9}
\bibitem{Dependencies} Project repository,
\emph{Arithmetic dependencies of the defining-polynomial constructions},
\nolinkurl{arithmetic_dependencies_20261108_v8.tex}, Section~3.
\end{thebibliography}
\end{document}
